\documentclass[conference]{IEEEtran}
\IEEEoverridecommandlockouts

\usepackage[spanish,es-tabla,es-nodecimaldot]{babel}
\usepackage[utf8]{inputenc}
\usepackage[T1]{fontenc}
\usepackage{mathptmx}
\usepackage{cite}
\usepackage{amsmath}
\usepackage{graphicx}
\usepackage{textcomp}
\usepackage{xcolor}
\usepackage{booktabs}
\usepackage{array}
\usepackage{siunitx}
\sisetup{output-decimal-marker={.},group-separator={\,},group-minimum-digits=4,detect-all}
\usepackage{url}
\usepackage{xurl}

% Marca visible para lo que falta (se elimina en la versión final)
\newcommand{\pendiente}[1]{\textcolor{red}{\textbf{[PENDIENTE: #1]}}}

\begin{document}

\title{Tecnologías de baterías y selección del banco de almacenamiento para un sistema fotovoltaico aislado}

\author{\IEEEauthorblockN{Abner López}
\IEEEauthorblockA{\textit{Ingeniería Electrónica} \\
\textit{Tecnológico de Costa Rica}\\
Cartago, Costa Rica}
\and
\IEEEauthorblockN{Sebastián Rojas}
\IEEEauthorblockA{\textit{Ingeniería Electrónica} \\
\textit{Tecnológico de Costa Rica}\\
Cartago, Costa Rica}
\and
\IEEEauthorblockN{Jordi Segura}
\IEEEauthorblockA{\textit{Ingeniería Electrónica} \\
\textit{Tecnológico de Costa Rica}\\
Cartago, Costa Rica}
}

\maketitle

\begin{abstract}
El banco de baterías es el componente que permite a un sistema fotovoltaico aislado alimentar las cargas durante la noche y en los días sin sol. En este informe se describen las tecnologías de baterías aplicables a un banco residencial de \SI{24}{V}: plomo-ácido de electrolito inundado, plomo-ácido regulada por válvula con separador de vidrio absorbente (AGM) y con electrolito en gel, y ion-litio de fosfato de hierro y litio (LiFePO\textsubscript{4}). Se definen los parámetros de comparación: capacidad nominal, tensión, profundidad de descarga, ciclos de vida, eficiencia, autodescarga, temperatura de operación, mantenimiento y costo. A partir de esta caracterización se identifican opciones comerciales con su ficha técnica y precio, se comparan mediante el costo por kWh útil y se selecciona la batería para el Proyecto~2. \pendiente{completar el resumen con la batería seleccionada y la capacidad del banco cuando se tenga el consumo}
\end{abstract}

\begin{IEEEkeywords}
baterías, sistema fotovoltaico aislado, plomo-ácido, AGM, gel, LiFePO\textsubscript{4}, profundidad de descarga, días de autonomía, análisis comparativo
\end{IEEEkeywords}

\section{Introducción}

En un sistema fotovoltaico conectado a la red, la energía que no se produce en un momento dado se toma de la red. En un sistema aislado esa opción no existe: si el sistema se dimensiona mal, las cargas se quedan sin energía. Por eso el banco de baterías es el componente que define la autonomía del sistema. Su capacidad debe cubrir el consumo nocturno y los días de baja radiación, y su tecnología determina cuánto de esa capacidad se puede usar, cuántos ciclos resiste antes de degradarse y cuánto cuesta cada kilovatio-hora que entrega.

El Proyecto~2 consiste en diseñar un sistema fotovoltaico aislado para la casa de Abner López, en Santa Rosa de Oreamuno, Cartago. Para el dimensionamiento se usa un banco de \SI{24}{V} y una profundidad de descarga del 70\,\%, valores fijados en clase como punto de partida. La casa tiene servicio eléctrico, por lo que el sistema se plantea como respaldo de un grupo de cargas críticas.

Este informe describe las tecnologías de baterías que se consideran para el proyecto y los parámetros con que se comparan. Luego identifica opciones comerciales de cada tecnología con su ficha técnica y su precio, las compara en una matriz y selecciona la batería para el proyecto. Finalmente, calcula la capacidad del banco a partir del consumo de las cargas críticas de la casa.

\section{Parámetros de comparación}\label{sec:parametros}

Las baterías se comparan con los parámetros que aparecen en sus dentro de sus hojas de datos y con base a sus especificaciones se dimenciona el banco. La Tabla~\ref{tab:parametros} resume qué indica cada uno y cómo se usa en el diseño.

La \emph{capacidad nominal} es la carga energetica que la batería entrega al descargarse a una corriente constante hasta un valor de  tensión final, en condiciones definidas por el fabricante \cite{linden}. Se expresa en Ah o, multiplicada por la tensión nominal, en Wh. Como la capacidad de la bateria disminuye cuando la corriente de descarga aumenta, el fabricante la declara para un tiempo de descarga: $C_{20}$ es la capacidad cuando la batería se descarga en \SI{20}{h}, y $C_{100}$ cuando se descarga en \SI{100}{h} \cite{linden}. En sistemas fotovoltaicos se usan normalmente $C_{20}$ y $C_{100}$, porque la descarga ocurre a lo largo de la noche con corrientes bajas. Para comparar modelos entre sí se usa la misma tasa en todos.

La \emph{profundidad de descarga} (DOD) es la fracción de la capacidad que se extrae en un ciclo, y el \emph{estado de carga} (SOC) es su complemento, $\mathrm{SOC} = 1 - \mathrm{DOD}$. La DOD se encarga de fijar qué parte de la capacidad nominal se puede usar y, a la vez, cuántos ciclos resiste la batería: entre más profunda es la descarga, menor es el número de ciclos \cite{linden}. Por eso los \emph{ciclos de vida} solo se pueden comparar si están declarados a la misma DOD.

La \emph{eficiencia} de ida y vuelta es la relación que se tiene entre la energía que la batería entrega y la que recibió para cargarse. La \emph{autodescarga} es la pérdida de carga de una batería que se encuentra en reposo, y se expresa como porcentaje por mes. La \emph{temperatura} corresponde a un parametro que afecta la capacidad disponible y la vida útil, y cada tecnología tiene un rango distinto para cargar y para descargar. El \emph{mantenimiento} indica si la batería depende de reposición de agua, cargas de ecualización o ventilación cada cierto tiempo. Además, la corriente máxima de carga y de descarga limita la potencia que el banco puede recibir del regulador y entregar al inversor, y puede obligar a conectar baterías en paralelo aunque la capacidad ya sea suficiente.

El \emph{costo} se compara por kWh útil, es decir, por la energía que se puede extraer en un ciclo sin superar la DOD recomendada:
\begin{equation}
c_{\text{útil}} = \frac{1000 \times \text{precio}}{V_{\mathrm{nom}}\,C\,\mathrm{DOD}},
\label{eq:costo_util}
\end{equation}
en USD/kWh, donde $V_{\mathrm{nom}}$ es la tensión nominal en V y $C$ la capacidad en Ah. Una batería barata por Ah puede resultar cara por kWh útil si su DOD recomendada es baja.

\begin{table}[!ht]
\caption{Parámetros de comparación de las baterías}
\label{tab:parametros}
\centering
\footnotesize
\setlength{\tabcolsep}{3pt}
\begin{tabular}{>{\raggedright\arraybackslash}p{0.27\columnwidth} >{\raggedright\arraybackslash}p{0.66\columnwidth}}
\toprule
Parámetro & Uso en el diseño \\
\midrule
Capacidad nominal (Ah, Wh) & Número de baterías en paralelo; se compara a la misma tasa de ($C_{20}$ o $C_{100}$). \\
Tensión nominal (V) & Número de baterías en serie para formar \SI{24}{V}. \\
DOD recomendada (\%) & Fracción utilizable de la capacidad; costo por kWh útil. \\
Ciclos de vida & Años hasta el reemplazo, a la DOD de diseño. \\
Eficiencia (\%) & Pérdidas en la batería dentro del balance de energía. \\
Autodescarga (\%/mes) & Pérdida en reposo; factor $K_a$ del balance. \\
Corriente máx. de carga y descarga (A) & Baterías en paralelo según el regulador y el inversor. \\
Temperatura de operación (°C) & Capacidad disponible y condiciones del lugar de instalación. \\
Mantenimiento & Agua, ecualización, ventilación, sistema de gestión (BMS). \\
Peso (kg) y garantía & Instalación en la casa y vida útil asegurada. \\
Costo (USD/kWh útil) & Comparación económica, ecuación~\eqref{eq:costo_util}. \\
\bottomrule
\end{tabular}
\end{table}

\section{Tecnologías de baterías}\label{sec:tecnologias}

Una batería recargable almacena energía mediante una reacción de reducción-oxidación reversible. Durante la descarga, el electrodo negativo se oxida y libera electrones que circulan por el circuito externo hasta el electrodo positivo, que se reduce; el electrolito permite el paso de iones entre ambos electrodos. Durante la carga, una fuente externa invierte la reacción \cite{linden}. El proceso no es completamente reversible: en cada ciclo parte del material activo deja de participar en la reacción, y por eso la capacidad disminuye con el uso.

Entre las tecnologías de almacenamiento electroquímico estudiadas se encuentran las baterías de plomo-ácido, ion-litio, níquel-cadmio, sodio-azufre y de flujo. Las tres últimas se descartan para el proyecto. Las de níquel-cadmio usan cadmio, un metal tóxico, y su costo instalado es de 1200 a \SI{1500}{USD/kWh}, el doble o más que el del plomo-ácido \cite{may2018}. Las de sodio-azufre operan a más de \SI{300}{\celsius} y las de flujo almacenan el electrolito en tanques externos; ambas están orientadas a instalaciones de gran escala de la red eléctrica \cite{may2018}, no a un banco residencial de \SI{24}{V}. El análisis se concentra entonces en el plomo-ácido, con sus tres variantes, y en el ion-litio de fosfato de hierro.

\subsection{Plomo-ácido}

La batería de plomo-ácido usa dióxido de plomo (PbO\textsubscript{2}) en el electrodo positivo, plomo esponjoso (Pb) en el negativo y una solución de ácido sulfúrico como electrolito. Durante la descarga ambos electrodos se convierten en sulfato de plomo y el ácido se consume, según
\begin{equation}
\mathrm{Pb + PbO_2 + 2H_2SO_4 \rightleftharpoons 2PbSO_4 + 2H_2O}.
\end{equation}
La tensión nominal de cada celda es de \SI{2}{V}; una batería de \SI{12}{V} tiene seis celdas en serie y un banco de \SI{24}{V} necesita doce \cite{linden}, \cite{pavlov}.

Es la tecnología recargable más madura. Se usa en vehículos, sistemas de respaldo y telecomunicaciones, su costo de batería es de 200 a \SI{500}{USD/kWh} y se recicla casi en su totalidad \cite{may2018}. Sus desventajas son el peso y la vida en ciclos profundos. Su densidad de energía es de 35 a \SI{40}{Wh/kg}, frente a 150 a \SI{180}{Wh/kg} del ion-litio \cite{may2018}. Si la batería permanece descargada, el sulfato de plomo forma cristales grandes que ya no se reconvierten durante la carga; este proceso, llamado sulfatación, reduce la capacidad de forma permanente \cite{pavlov}. Por esa razón la descarga se limita: en el curso se usa una DOD máxima del 70\,\%, y muchos fabricantes declaran sus ciclos al 50\,\%. La práctica recomendada IEEE~1013 trata el dimensionamiento de este tipo de baterías en sistemas fotovoltaicos aislados \cite{ieee1013}.

Las variantes de plomo-ácido se diferencian en cómo se contiene el electrolito.

\subsubsection{Electrolito inundado}

En la batería inundada el electrolito es líquido y cubre libremente las placas. Al final de la carga parte de la corriente descompone el agua en hidrógeno y oxígeno, que salen por los tapones de ventilación. Esto obliga a reponer agua destilada de forma periódica, a instalar la batería en posición vertical y a ventilar el lugar para que el hidrógeno no se acumule \cite{linden}. Además, el ácido tiende a concentrarse en la parte inferior de la celda, por lo que el regulador debe aplicar cargas de ecualización: una sobrecarga controlada que produce gas y mezcla el electrolito \cite{pavlov}.

A cambio de ese mantenimiento, las baterías inundadas de ciclo profundo, como las de placas gruesas para carritos de golf o las de placa tubular OPzS, tienen el menor costo por Ah de la familia y toleran bien las sobrecargas, porque el agua perdida se puede reponer \cite{pavlov}. Para una vivienda, el mantenimiento y la necesidad de un espacio ventilado son las principales limitaciones.

\subsubsection{Regulada por válvula con separador AGM}

Las baterías reguladas por válvula (VRLA) son selladas. El oxígeno que se produce en el electrodo positivo al final de la carga se difunde hasta el negativo y se recombina para formar agua otra vez, de modo que la pérdida de agua es mínima y no hay que reponerla. Una válvula se abre solo si la presión interna supera un límite \cite{linden}, \cite{pavlov}.

En la variante AGM (\emph{absorbent glass mat}), el electrolito está absorbido en un separador de microfibra de vidrio colocado entre las placas. El separador no está saturado del todo, y los poros libres son los que permiten el paso del oxígeno hacia el electrodo negativo \cite{pavlov}. Como el electrolito queda retenido y muy cerca de las placas, la resistencia interna es baja y la batería entrega corrientes altas. No requiere reposición de agua y se puede instalar en otras posiciones, pero es sensible a la sobrecarga y a la temperatura alta: si la tensión de carga es excesiva, el gas que se libera por la válvula ya no se recupera y la batería se seca \cite{linden}. Por eso el regulador debe configurarse con las tensiones de carga que indica el fabricante.

\subsubsection{Regulada por válvula con electrolito en gel}

En la batería de gel, el ácido sulfúrico se mezcla con sílice pirogénica y forma un gel que inmoviliza el electrolito. Con el tiempo el gel se agrieta, y esas grietas son las que permiten la recombinación del oxígeno \cite{pavlov}. Igual que la AGM, es sellada y no requiere reposición de agua.

El gel evita que el ácido se estratifique y protege mejor a las placas en descargas profundas, por lo que estas baterías resisten más ciclos que las AGM en aplicaciones de ciclo diario, como las solares \cite{pavlov}. A cambio, su resistencia interna es mayor, entregan menos corriente pico y admiten una corriente de carga menor. Las versiones de placa tubular, conocidas como OPzV, son las de mayor vida en ciclos dentro del plomo-ácido sellado \cite{pavlov}.

\subsection{Ion-litio de fosfato de hierro (LiFePO\textsubscript{4})}

En una batería de ion-litio, los iones de litio se desplazan entre dos electrodos que los alojan dentro de su estructura cristalina: durante la carga salen del electrodo positivo y se insertan en el negativo, de grafito, y durante la descarga hacen el recorrido inverso \cite{tarascon2001}. El electrolito es una sal de litio disuelta en un solvente orgánico. Como el litio no se deposita como metal en ningún momento del ciclo normal, la reacción es muy reversible y la eficiencia de ida y vuelta es alta.

La química del electrodo positivo define el tipo de batería. Las de óxido de cobalto o de níquel-manganeso-cobalto (NMC) tienen mayor densidad de energía y se usan en electrónica portátil y vehículos eléctricos. La de fosfato de hierro y litio, LiFePO\textsubscript{4} o LFP, fue propuesta por Padhi \emph{et al.}\ en 1997 \cite{padhi1997}. En su estructura de olivino, el oxígeno está unido con fuerza al fósforo, lo que hace más difícil que se libere al calentarse; por eso las celdas LFP son más estables térmicamente que las de óxidos y tienen menor riesgo de incendio \cite{zubi2018}. Además, no usan cobalto ni níquel. La tensión nominal de cada celda es de \SI{3.2}{V}, de modo que ocho celdas en serie forman una batería de \SI{25.6}{V}, que es la versión comercial de \SI{24}{V} \cite{linden}.

Para el uso en sistemas aislados, las ventajas de la LFP son la profundidad de descarga y la vida en ciclos. Los fabricantes declaran ciclos con descargas del 80 al 100\,\%, frente al 50\,\% habitual en el plomo-ácido, y su densidad de energía es del orden de cuatro veces la del plomo-ácido \cite{may2018}, \cite{zubi2018}. No requiere mantenimiento ni ecualización, y su autodescarga es menor que la del plomo-ácido \cite{linden}.

Sus limitaciones son tres. Primero, cada batería necesita un sistema de gestión (BMS, \emph{battery management system}) que equilibre la tensión de las celdas y la desconecte por sobretensión, subtensión, sobrecorriente o temperatura fuera de rango \cite{linden}. Segundo, no se debe cargar por debajo de \SI{0}{\celsius}, porque el litio se deposita como metal sobre el grafito y daña la celda \cite{linden}. Tercero, la curva de descarga es muy plana, por lo que el estado de carga no se puede estimar solo con la tensión y se requiere un monitor que cuente la carga entrante y saliente. En costo, May \emph{et al.} reportan de 340 a \SI{590}{USD/kWh} para baterías LFP de la red eléctrica \cite{may2018}; los precios de baterías comerciales de \SI{24}{V} se revisan en la Sección~\ref{sec:opciones}.

\subsection{Síntesis}

La Tabla~\ref{tab:tecnologias} resume las diferencias cualitativas entre las cuatro tecnologías. Los valores numéricos de cada modelo comercial, tomados de sus fichas técnicas, se comparan en la Sección~\ref{sec:matriz}. Para todas ellas, la norma IEC~61427-1 define los ensayos de baterías destinadas a sistemas fotovoltaicos aislados \cite{iec61427}.

\begin{table}[!t]
\caption{Síntesis de las tecnologías consideradas}
\label{tab:tecnologias}
\centering
\footnotesize
\setlength{\tabcolsep}{2.5pt}
\begin{tabular}{>{\raggedright\arraybackslash}p{0.21\columnwidth} >{\raggedright\arraybackslash}p{0.17\columnwidth} >{\raggedright\arraybackslash}p{0.17\columnwidth} >{\raggedright\arraybackslash}p{0.17\columnwidth} >{\raggedright\arraybackslash}p{0.18\columnwidth}}
\toprule
 & Inundada & AGM & Gel & LiFePO\textsubscript{4} \\
\midrule
Tensión de celda & \SI{2}{V} & \SI{2}{V} & \SI{2}{V} & \SI{3.2}{V} \\
Celdas para \SI{24}{V} & 12 & 12 & 12 & 8 \\
Electrolito & Líquido libre & Absorbido en fibra de vidrio & Gelificado con sílice & Sal de litio en solvente orgánico \\
Reposición de agua & Sí & No & No & No \\
Ecualización & Sí & No & No & No \\
Ventilación & Necesaria & Reducida & Reducida & No \\
Requiere BMS & No & No & No & Sí \\
Principal limitación & Manteni\-miento & Sensible a sobrecarga & Corriente pico menor & Costo inicial; no cargar bajo \SI{0}{\celsius} \\
\bottomrule
\end{tabular}
\end{table}

\section{Opciones comerciales}\label{sec:opciones}

La búsqueda se orientó a modelos de ciclo profundo que permitan formar un banco de \SI{24}{V}. Se eligieron cinco modelos de dos tecnologías del curso: uno de cada variante de plomo-ácido y dos de LiFePO\textsubscript{4}, uno de bajo costo y otro de un fabricante de equipos para sistemas aislados. Los precios se consultaron el 5 de noviembre de 2026. Se priorizaron tiendas de Costa Rica, pero solo se encontró una referencia local, para B4, y sin el modelo exacto. Por eso se usan precios minoristas de Estados Unidos y, cuando no los hubo, de Europa. Los euros se convirtieron con 1~USD = 0.8879~EUR \cite{eurusd}. La Tabla~\ref{tab:opciones} resume las especificaciones.

\subsection{B1: Trojan T-105}

Batería de plomo-ácido inundada de \SI{6}{V}, de placas gruesas para ciclo profundo. Su capacidad es de \SI{225}{Ah} a $C_{20}$ y \SI{250}{Ah} a $C_{100}$, con 1200 ciclos al 50\,\% de DOD y una masa cercana a \SI{28}{kg} \cite{trojan_sig}. La autodescarga declarada es de 5 a 15\,\% por mes, según la temperatura de almacenamiento \cite{trojan_t105}. Requiere reposición de agua y cargas de ecualización. Para \SI{24}{V} se necesitan cuatro en serie. Precio: \SI{225.00}{USD} por unidad \cite{res_t105}.

\subsection{B2: Trojan SAGM 12 205}

Batería de plomo-ácido regulada por válvula, con separador AGM, de la línea solar de Trojan. Tiene \SI{12}{V} y \SI{205}{Ah} a $C_{20}$, 1700 ciclos al 50\,\% de DOD, una masa de \SI{59}{kg} y dimensiones de $380 \times 176 \times \SI{357}{mm}$ \cite{trojan_sagm}. Para \SI{24}{V} se necesitan dos en serie. Precio: \SI{799.95}{USD} por unidad \cite{air1_sagm}.

\subsection{B3: Ultracell UCG200-12}

Batería de plomo-ácido regulada por válvula, con electrolito en gel, de \SI{12}{V} y \SI{200}{Ah}. Resiste unos 1200 ciclos al 50\,\% de DOD y 2400 al 30\,\%. Puede almacenarse hasta nueve meses a \SI{25}{\celsius} antes de requerir una carga de refresco, conserva el 86\,\% de su capacidad a \SI{0}{\celsius} y su masa es de \SI{62.9}{kg} \cite{ultracell}. Para \SI{24}{V} se necesitan dos en serie. Precio: 358.38~EUR con IVA europeo incluido, equivalentes a \SI{403.63}{USD} \cite{wattuneed}.

\subsection{B4: Power Queen LiFePO\textsubscript{4} 24V 100Ah}

Batería de LiFePO\textsubscript{4} de \SI{25.6}{V} y \SI{100}{Ah} (\SI{2.56}{kWh}), formada por ocho celdas en serie, con un BMS integrado de \SI{100}{A} de descarga continua y \SI{280}{A} durante \SI{5}{s}. El fabricante declara más de 4000 ciclos y una garantía de cinco años, y su masa es de unos \SI{20.7}{kg} \cite{pq}. La documentación encontrada no indica a qué DOD corresponden los 4000 ciclos. Una sola unidad forma el banco de \SI{24}{V}. Precio: \SI{399.99}{USD} \cite{pq_walmart}. En Costa Rica, el distribuidor oficial de la marca anuncia baterías de litio de \SI{24}{V} desde 400\,000 colones \cite{kuat}, sin indicar el modelo.

\subsection{B5: Victron LiFePO\textsubscript{4} 25.6V/200Ah Smart-a}

Batería de LiFePO\textsubscript{4} de \SI{25.6}{V} y \SI{200}{Ah} (\SI{5.12}{kWh}), con balanceo de celdas y monitoreo por Bluetooth. Resiste 2500 ciclos al 80\,\% de DOD, 3000 al 70\,\% y 5000 al 50\,\%, y su masa es de \SI{39}{kg} \cite{victron_200}. La eficiencia de ida y vuelta declarada para la línea es de 92\,\% \cite{victron_ds}. Una unidad forma el banco de \SI{24}{V}. Precio: \SI{2983.50}{USD} \cite{victron_precio}.

\begin{table}[!t]
\caption{Especificaciones de las opciones comerciales}
\label{tab:opciones}
\centering
\footnotesize
\setlength{\tabcolsep}{2.2pt}
\begin{tabular}{l l c c c c c}
\toprule
Ref. & Tecnología & $V$ [V] & $C$ [Ah] & Unid. 24\,V & Ciclos (DOD) & Masa [kg] \\
\midrule
B1 & Inundada & 6 & 225 & 4 & 1200 (50\,\%) & 28 \\
B2 & AGM & 12 & 205 & 2 & 1700 (50\,\%) & 59 \\
B3 & Gel & 12 & 200 & 2 & 1200 (50\,\%) & 62.9 \\
B4 & LiFePO\textsubscript{4} & 25.6 & 100 & 1 & $>$4000 (n.d.) & 20.7 \\
B5 & LiFePO\textsubscript{4} & 25.6 & 200 & 1 & 2500 (80\,\%) & 39 \\
\bottomrule
\end{tabular}
\\[2pt]
{\footnotesize Capacidad a $C_{20}$ para B1 a B3. n.d.: no declarado en la documentación consultada.}
\end{table}

\section{Matriz comparativa}\label{sec:matriz}

La Tabla~\ref{tab:matriz} compara las cinco opciones con un módulo de \SI{24}{V} armado con cada modelo, es decir, con las unidades en serie de la Tabla~\ref{tab:opciones}. El costo por kWh útil se calculó con la ecuación~\eqref{eq:costo_util}, usando la DOD a la que cada fabricante declara sus ciclos: 50\,\% para el plomo-ácido y 80\,\% para el LiFePO\textsubscript{4}. También se incluye el costo con la DOD de diseño del curso, 70\,\%, y el costo por kWh entregado durante toda la vida en ciclos,
\begin{equation}
c_{\text{vida}} = \frac{c_{\text{útil}}}{N_{\text{ciclos}}},
\label{eq:costo_vida}
\end{equation}
que reúne en una cifra el precio, la DOD y los ciclos. Los cálculos están en \texttt{calculos/baterias.py}.

\begin{table*}[!t]
\caption{Matriz comparativa de las baterías (módulo de \SI{24}{V})}
\label{tab:matriz}
\centering
\footnotesize
\setlength{\tabcolsep}{4pt}
\begin{tabular}{l r c c r r c r c >{\raggedright\arraybackslash}p{0.17\textwidth} r c}
\toprule
Ref. & Precio & $E_{\text{nom}}$ & DOD & \multicolumn{2}{c}{Costo por kWh útil [USD]} & Ciclos & $c_{\text{vida}}$ & Eficiencia & Mantenimiento & Masa & Garantía \\
\cmidrule(lr){5-6}
 & [USD] & [kWh] & tecnol. & DOD tecnol. & DOD 70\,\% & & [USD/kWh] & [\%] & & [kg] & [años] \\
\midrule
B1 & 900.00 & 5.40 & 50\,\% & 333 & 238 & 1200 & 0.28 & n.d. & Agua, ecualización, ventilación & 112 & n.c.\textsuperscript{a} \\
B2 & 1599.90 & 4.92 & 50\,\% & 650 & 465 & 1700 & 0.38 & n.d. & Ninguno; tensiones de carga AGM & 118 & 2\textsuperscript{b} \\
B3 & 807.25 & 4.80 & 50\,\% & 336 & 240 & 1200 & 0.28 & n.d. & Ninguno; tensiones de carga gel & 126 & n.d. \\
B4 & 399.99 & 2.56 & 80\,\% & 195 & 223 & $>$4000\textsuperscript{c} & 0.05\textsuperscript{c} & n.d. & Ninguno; BMS integrado & 21 & 5 \\
B5 & 2983.50 & 5.12 & 80\,\% & 728 & 832 & 2500 & 0.29 & 92 & Ninguno; confirmar si requiere BMS externo & 39 & 5\textsuperscript{d} \\
\bottomrule
\end{tabular}
\\[2pt]
{\footnotesize Precio y masa del conjunto de unidades que forma \SI{24}{V}. n.d.: no declarado en la documentación consultada. n.c.: datos contradictorios.
\textsuperscript{a}Los distribuidores indican 1 o 5 años. \textsuperscript{b}Según distribuidor \cite{solaris_sagm}. \textsuperscript{c}La DOD de los 4000 ciclos no está declarada; $c_{\text{vida}}$ se calculó suponiendo que corresponden al 80\,\%. \textsuperscript{d}Valor más repetido entre distribuidores; algunos indican 3 años.}
\end{table*}

\subsection{Costo por kWh útil}

Por precio total, B4 es la opción más barata, con \SI{399.99}{USD}, y B5 la más cara, con \SI{2983.50}{USD}. Sin embargo, B4 almacena \SI{2.56}{kWh}, mientras que los módulos de plomo-ácido almacenan entre 4.80 y \SI{5.40}{kWh}, por lo que el precio total no basta para comparar. Por kWh útil, B4 cuesta \SI{195}{USD}, B1 y B3 cuestan 333 y \SI{336}{USD}, y B2 y B5 cuestan 650 y \SI{728}{USD}. La diferencia entre B4 y el plomo-ácido viene de la DOD: el plomo-ácido solo entrega la mitad de su capacidad nominal sin acortar su vida, y el LiFePO\textsubscript{4} entrega el 80\,\%.

Con la DOD de diseño del curso, 70\,\%, B1, B3 y B4 quedan entre 223 y \SI{240}{USD/kWh} útil. Esa comparación favorece al plomo-ácido, porque descargarlo al 70\,\% le resta ciclos respecto a los declarados al 50\,\%.

\subsection{Ciclos y costo durante la vida útil}

Los modelos de plomo-ácido declaran entre 1200 y 1700 ciclos al 50\,\% de DOD. B5 declara 2500 ciclos al 80\,\%, y B4 más de 4000 sin indicar la DOD. Con un ciclo por día, 1200 ciclos equivalen a unos 3.3 años y 2500 ciclos a unos 6.8 años. El costo por kWh entregado en la vida útil, ecuación~\eqref{eq:costo_vida}, es de 0.28 a \SI{0.38}{USD/kWh} para el plomo-ácido y de \SI{0.29}{USD/kWh} para B5. Para B4 resulta de \SI{0.05}{USD/kWh} si los 4000 ciclos corresponden al 80\,\% de DOD; esta cifra es la más incierta de la matriz y debe confirmarse con la ficha oficial.

\subsection{Eficiencia y autodescarga}

Solo B5 declara la eficiencia de ida y vuelta, 92\,\% \cite{victron_ds}. Para las demás se usan como referencia los valores por tecnología de May \emph{et al.}: 85\,\% para el plomo-ácido y 90\,\% para el ion-litio \cite{may2018}. Con esos valores, por cada kWh que entregan, las baterías de plomo-ácido requieren unos \SI{1.18}{kWh} de carga, frente a \SI{1.09}{kWh} de B5. En autodescarga, B1 pierde de 5 a 15\,\% por mes \cite{trojan_t105}, mientras que B3 puede almacenarse nueve meses antes de recargarla \cite{ultracell}. En un sistema que se carga todos los días, la autodescarga está cubierta por el factor $K_a$ del balance de energía.

\subsection{Mantenimiento, peso y garantía}

B1 es la única que requiere reposición de agua, ecualización y un lugar ventilado. Las demás no tienen mantenimiento periódico, pero B2 y B3 exigen que el regulador use las tensiones de carga de su tecnología. El banco de plomo-ácido de \SI{24}{V} pesa entre 112 y \SI{126}{kg} en dos o cuatro unidades de 28 a \SI{63}{kg} cada una. B4 pesa \SI{21}{kg} y B5 \SI{39}{kg}, con un solo bloque. En garantía, B4 y B5 ofrecen cinco años, B2 dos años, y para B1 y B3 los datos son contradictorios o no aparecen.

\section{Dimensionamiento del banco}\label{sec:dimensionamiento}

\pendiente{procedimiento del curso ($Q = 1.1\,E_T\,D/(V\,\mathrm{DOD})$, \SI{24}{V}, DOD de 70\,\%) y tabla provisional de capacidad para varios consumos diarios; el consumo real de la casa está pendiente.}

\section{Selección preliminar}\label{sec:seleccion}

\pendiente{batería seleccionada y justificación; la cantidad y la configuración del banco se confirman con el consumo.}


\begin{thebibliography}{99}

\bibitem{linden} T.~B. Reddy, Ed., \emph{Linden's Handbook of Batteries}, 4th~ed. New York, NY, USA: McGraw-Hill, 2011.

\bibitem{pavlov} D.~Pavlov, \emph{Lead-Acid Batteries: Science and Technology}, 2nd~ed. Amsterdam, The Netherlands: Elsevier, 2017.

\bibitem{may2018} G.~J. May, A.~Davidson y B.~Monahov, ``Lead batteries for utility energy storage: A review,'' \emph{J. Energy Storage}, vol.~15, pp.~145--157, 2018.

\bibitem{tarascon2001} J.-M. Tarascon y M.~Armand, ``Issues and challenges facing rechargeable lithium batteries,'' \emph{Nature}, vol.~414, pp.~359--367, 2001.

\bibitem{padhi1997} A.~K. Padhi, K.~S. Nanjundaswamy y J.~B. Goodenough, ``Phospho-olivines as positive-electrode materials for rechargeable lithium batteries,'' \emph{J. Electrochem. Soc.}, vol.~144, no.~4, pp.~1188--1194, 1997.

\bibitem{zubi2018} G.~Zubi, R.~Dufo-López, M.~Carvalho y G.~Pasaoglu, ``The lithium-ion battery: State of the art and future perspectives,'' \emph{Renew. Sustain. Energy Rev.}, vol.~89, pp.~292--308, 2018.

\bibitem{ieee1013} \emph{IEEE Recommended Practice for Sizing Lead-Acid Batteries for Stand-Alone Photovoltaic (PV) Systems}, IEEE Std 1013-2019, 2019.

\bibitem{iec61427} \emph{Secondary cells and batteries for renewable energy storage -- General requirements and methods of test -- Part~1: Photovoltaic off-grid application}, IEC 61427-1:2013, 2013.

\end{thebibliography}

\end{document}
